% Generated draft: requires mathematical and source review.
\documentclass[11pt]{article}
\usepackage[margin=25mm]{geometry}
\usepackage{fontspec}
\setmainfont{Latin Modern Roman}
\usepackage{amsmath,amssymb,mathtools,unicode-math}
\setmathfont{Latin Modern Math}
\usepackage{enumitem,needspace,xurl,hyperref}
\setlength{\parindent}{0pt}
\setlength{\parskip}{5pt}
\setlength{\emergencystretch}{3em}
\newcommand{\N}{\mathbb{N}}
\newcommand{\Z}{\mathbb{Z}}
\newcommand{\Q}{\mathbb{Q}}
\newcommand{\R}{\mathbb{R}}
\newcommand{\C}{\mathbb{C}}
\newcommand{\F}{\mathbb{F}}
\newcommand{\A}{\mathbb{A}}
\newcommand{\PP}{\mathbb{P}}
\newcommand{\E}{\mathbb{E}}
\newcommand{\eps}{\varepsilon}
\newcommand{\dd}{\,\mathrm{d}}
\begin{document}
\section*{UnsolvedMath ID: 701325}
\textbf{Source problem code:} catalogue entry
\subsection{Definitions and mathematical statement}
Let \(\Omega\) be a countably infinite set, and let \(S=\operatorname{Sym}(\Omega)\) be the group of all bijections \(\Omega\to\Omega\), with composition as the group operation. For \(A\subseteq S\), let \(\langle A\rangle\) denote the subgroup generated algebraically by \(A\), consisting of finite products of elements of \(A\) and their inverses. Say that \(S\) is finitely generated over a subgroup \(K\leq S\) if there exists a finite set \(F\subseteq S\) such that \(S=\langle K\cup F\rangle\).

Is it true that, for every pair of subgroups \(G,H\leq S\) satisfying \(\langle G\cup H\rangle=S\), one has
\[
\bigl(\exists\,F_G\subseteq S\text{ finite}:\ S=\langle G\cup F_G\rangle\bigr)
\quad\text{or}\quad
\bigl(\exists\,F_H\subseteq S\text{ finite}:\ S=\langle H\cup F_H\rangle\bigr)?
\]
Here \(S\) is the full symmetric group, including permutations of infinite support; generation means algebraic generation, with no topological closure involved.
\subsection{Short English statement}
If two subgroups together generate the full symmetric group on a countably infinite set, must adjoining finitely many permutations to at least one of those subgroups generate the whole symmetric group?
\subsection{Sources}
\begin{itemize}
\item Local UnsolvedMath registry, supplied record with permanent ID \(701325\), exact source problem code \texttt{catalogue entry}, titled ``Kourovka Notebook Problem 17.23.''
\item The supplied record attributes Problem 17.23 to G. M. Bergman, proposed in 2010, in \emph{The Kourovka Notebook (Unsolved Problems in Group Theory)}, 17th issue (2010). It also identifies the 21st issue edition (2026), arXiv:1401.0300v48. The registry describes the collection as edited by E. I. Khukhro and V. D. Mazurov. These bibliographic details are supplied source information and have not been independently verified.
\end{itemize}
\end{document}
